\chapter*{Wstęp}
\addcontentsline{toc}{chapter}{Wstęp}

Krajowy System e-Faktur, w skrócie KSeF, to platforma online pozwalająca na wystawianie, odbieranie oraz przechowywanie faktur ustrukturyzowanych, opracowana przez Ministerstwo Finansów \parencite{ministerstwo_finansow_pytania_2025}. 

Celem pracy jest zaprojektowanie oraz utworzenie oprogramowania umożliwiającego bezpieczną, nieinwazyjną i minimalizującą koszty integrację z Krajowym Systemem e-Faktur, w szczególności dla małych i średnich przedsiębiorstw. Projekt zrealizowany został w języku Go, z wykorzystaniem innych technologii takich jak SQLite oraz HTMX i przybrał formę pojedynczego pliku wykonywalnego, kompatybilnego z systemami Windows, macOS oraz Linux. 
\par
Zakres pracy obejmuje:
\begin{itemize}
    \item analizę wymagań technicznych oraz oficjalnej dokumentacji Krajowego Systemu e-Faktur
    \item omówienie technologii wykorzystanych do utworzenia oprogramowania,
    \item zaprojektowanie architektury oprogramowania,
    \item implementację oprogramowania,
    \item testowanie jego działania,
    \item praktyczne wykorzystanie go w testowym scenariuszu oraz ocenę działania na podstawie wyników.
\end{itemize}

Niniejsza praca ustrukturyzowana jest według powyższej listy. W rozdziale pierwszym omówione są ogólne założenia projektu, środowisko Krajowego Systemu e-Faktur oraz ogólna architektura oprogramowania. W rozdziale drugim wytłumaczone są pojęcia wymagane do zrozumienia rozdziału trzeciego oraz omówione są technologie wykorzystane do realizacji projektu, takie jak język programowania czy formaty elektronicznej wymiany danych. Rozdział trzeci składa się z omówienia architektury utworzonej aplikacji, jej komponentów oraz funkcji, wraz z przedstawieniem przykładowych przeprowadzonych scenariuszy testowych.
